\documentclass[conference, onecolumn]{IEEEtran}
\usepackage[utf8]{inputenc}
\usepackage[portuguese]{babel}
\usepackage{graphicx}
\usepackage{float}
\usepackage{listings}
\usepackage{xcolor}
\usepackage{booktabs}
\usepackage{amsmath}
\usepackage{caption}
\usepackage{longtable}
\usepackage[hidelinks]{hyperref}

\lstset{
    basicstyle=\ttfamily\footnotesize,
    backgroundcolor=\color{gray!10},
    frame=single,
    breaklines=true,
    columns=fullflexible,
    keepspaces=true,
    tabsize=4,
    showstringspaces=false,
    numbers=left,
    numberstyle=\tiny\color{gray},
    stepnumber=1,
    numbersep=5pt
}

% =========================================================
\title{Trabalho Prático 1 --- Teste Funcional\\
\large Confiabilidade e Segurança de Software}

\author{
    \IEEEauthorblockN{Bruno Becker Silva}
    \IEEEauthorblockA{Escola Politécnica, PUCRS \\ \texttt{bruno.becker@edu.pucrs.br}}
}

\begin{document}
\maketitle

% =========================================================
\section*{Resumo}

Este relatório apresenta a suíte de testes funcionais de caixa-preta desenvolvida para a
função \texttt{calcular\_dano()}, aplicando as técnicas de Particionamento de Equivalência
(PE) e Análise do Valor Limite (BVA). Foram projetados 47 casos de teste que detectam
corretamente defeitos em 10 versões defeituosas da implementação.

% =========================================================
\section{Especificação da Função}

A função sob teste possui a seguinte assinatura em C:

\begin{lstlisting}[language=C, numbers=none]
int calcular_dano(int poder_ataque, int defesa, int tipo_ataque, int eh_critico);
\end{lstlisting}

\subsection{Domínio dos Parâmetros}

\begin{table}[H]
\centering
\caption{Domínio de cada parâmetro de \texttt{calcular\_dano}.}
\begin{tabular}{@{}llll@{}}
\toprule
\textbf{Parâmetro} & \textbf{Tipo} & \textbf{Domínio} & \textbf{Descrição} \\
\midrule
\texttt{poder\_ataque} & int & {[}1, 150{]} & Poder de ataque do personagem \\
\texttt{defesa}        & int & {[}0, 150{]} & Defesa do alvo \\
\texttt{tipo\_ataque}  & int & \{1, 2, 3\}  & 1=físico, 2=mágico, 3=verdadeiro \\
\texttt{eh\_critico}   & int & \{0, 1\}     & 0=normal, 1=golpe crítico \\
\bottomrule
\end{tabular}
\end{table}

\subsection{Regras de Cálculo (em ordem)}

\begin{enumerate}
    \item \textbf{Validação:} qualquer parâmetro fora do domínio $\Rightarrow$ retorna $-1$
    \item \textbf{Dano base:}
    \begin{itemize}
        \item Físico (1): $\text{dano} = \texttt{poder\_ataque} - \texttt{defesa}$
        \item Mágico (2): $\text{dano} = \texttt{poder\_ataque} - \lfloor\texttt{defesa}/2\rfloor$
        \item Verdadeiro (3): $\text{dano} = \texttt{poder\_ataque}$ (ignora \texttt{defesa})
    \end{itemize}
    \item \textbf{Dano mínimo:} se $\text{dano} < 1 \Rightarrow \text{dano} = 1$
    \item \textbf{Golpe crítico:} se $\texttt{eh\_critico} = 1 \Rightarrow \text{dano} = \text{dano} \times 2$
    \item \textbf{Retorno:} valor final de \texttt{dano}
\end{enumerate}

% =========================================================
\section{Entregável 1 — Classes de Equivalência}

\subsection{Classes Inválidas (CEI)}

\begin{table}[H]
\centering
\caption{Classes de equivalência inválidas.}
\begin{tabular}{@{}lllp{5.5cm}@{}}
\toprule
\textbf{Classe} & \textbf{Parâmetro} & \textbf{Condição} & \textbf{Resultado esperado} \\
\midrule
CEI-1 & \texttt{poder\_ataque} & $< 1$   & Retorna $-1$ (invalida toda a chamada) \\
CEI-2 & \texttt{poder\_ataque} & $> 150$ & Retorna $-1$ \\
CEI-3 & \texttt{defesa}        & $< 0$   & Retorna $-1$ \\
CEI-4 & \texttt{defesa}        & $> 150$ & Retorna $-1$ \\
CEI-5 & \texttt{tipo\_ataque}  & $< 1$   & Retorna $-1$ \\
CEI-6 & \texttt{tipo\_ataque}  & $> 3$   & Retorna $-1$ \\
CEI-7 & \texttt{eh\_critico}   & $< 0$   & Retorna $-1$ \\
CEI-8 & \texttt{eh\_critico}   & $> 1$   & Retorna $-1$ \\
\bottomrule
\end{tabular}
\end{table}

\subsection{Classes Válidas (CEV)}

\begin{table}[H]
\centering
\caption{Classes de equivalência válidas.}
\begin{tabular}{@{}llp{7cm}@{}}
\toprule
\textbf{Classe} & \textbf{Combinação} & \textbf{Descrição} \\
\midrule
CEV-1  & Físico / normal / dano$>$0        & \texttt{poder} $>$ \texttt{defesa}: dano base positivo sem crítico \\
CEV-2  & Físico / normal / dano$\leq$0     & \texttt{poder} $\leq$ \texttt{defesa}: dano mínimo ativado \\
CEV-3  & Físico / crítico / dano$>$0       & Dano base positivo com multiplicador crítico \\
CEV-4  & Físico / crítico / dano$\leq$0    & Dano mínimo ativado, depois crítico (ordem das regras 3 e 4) \\
CEV-5  & Mágico / normal / dano$>$0        & Dano base $= \texttt{poder} - \lfloor\texttt{defesa}/2\rfloor > 0$ \\
CEV-6  & Mágico / normal / dano$\leq$0     & Dano base negativo, mínimo ativado \\
CEV-7  & Mágico / crítico / dano$>$0       & Dano mágico positivo com multiplicador crítico \\
CEV-8  & Mágico / crítico / dano$\leq$0    & Mínimo ativado antes do crítico \\
CEV-9  & Verdadeiro / normal               & \texttt{defesa} ignorada; dano $= \texttt{poder}$ \\
CEV-10 & Verdadeiro / crítico              & \texttt{defesa} ignorada; dano $= \texttt{poder} \times 2$ \\
\bottomrule
\end{tabular}
\end{table}

% =========================================================
\section{Entregável 2 — Análise do Valor Limite (BVA)}

\begin{table}[H]
\centering
\caption{Valores on/off para cada parâmetro e para a fronteira do dano mínimo.}
\begin{tabular}{@{}lllll@{}}
\toprule
\textbf{Parâmetro} & \textbf{On-min} & \textbf{Off-min} & \textbf{On-max} & \textbf{Off-max} \\
\midrule
\texttt{poder\_ataque} & 1 (válido) & 0 $\to -1$ & 150 (válido) & 151 $\to -1$ \\
\texttt{defesa}        & 0 (válido) & $-1 \to -1$ & 150 (válido) & 151 $\to -1$ \\
\texttt{tipo\_ataque}  & 1 (válido) & 0 $\to -1$ & 3 (válido) & 4 $\to -1$ \\
\texttt{eh\_critico}   & 0 (válido) & $-1 \to -1$ & 1 (válido) & 2 $\to -1$ \\
\midrule
\multicolumn{5}{l}{\textbf{Fronteira do dano mínimo} (dano\_base $\in \{0, 1\}$, com e sem crítico)} \\
\midrule
dano\_base = 0 & \multicolumn{4}{l}{mínimo ativado $\Rightarrow$ retorna 1 (ou 2 com crítico)} \\
dano\_base = 1 & \multicolumn{4}{l}{mínimo \textit{não} ativado $\Rightarrow$ retorna 1 (mas por caminho diferente)} \\
dano\_base = 0 + crítico & \multicolumn{4}{l}{mínimo $\to$ 1 $\to$ crítico $\to$ 2 (testa ordem das regras 3 e 4)} \\
\bottomrule
\end{tabular}
\end{table}

% =========================================================
\section{Entregável 3 — Suíte de Testes}

A Tabela~\ref{tab:suite} apresenta os 47 casos de teste projetados, organizados em 8 blocos.

\begin{longtable}{@{}p{0.5cm}p{0.8cm}p{0.8cm}p{0.5cm}p{0.7cm}p{1.0cm}p{0.9cm}p{3.5cm}@{}}
\caption{Suíte completa de 47 casos de teste para \texttt{calcular\_dano}.}
\label{tab:suite} \\
\toprule
\textbf{ID} & \textbf{poder} & \textbf{defesa} & \textbf{tipo} & \textbf{crit.} & \textbf{Esperado} & \textbf{Técnica} & \textbf{Classe/Fronteira} \\
\midrule
\endfirsthead
\multicolumn{8}{l}{\small\textit{Tabela~\ref{tab:suite} — continuação}} \\
\toprule
\textbf{ID} & \textbf{poder} & \textbf{defesa} & \textbf{tipo} & \textbf{crit.} & \textbf{Esperado} & \textbf{Técnica} & \textbf{Classe/Fronteira} \\
\midrule
\endhead
\midrule
\multicolumn{8}{r}{\small\textit{continua na próxima página}} \\
\endfoot
\bottomrule
\endlastfoot
% --- Bloco 1: Classes Inválidas ---
\multicolumn{8}{l}{\textit{Bloco 1 — Classes Inválidas (PE)}} \\
T01 & 0    & 10  & 1 & 0 & $-1$ & PE & CEI-1: poder $< 1$ (zero) \\
T02 & $-5$ & 10  & 1 & 0 & $-1$ & PE & CEI-1: poder $< 1$ (negativo) \\
T03 & 151  & 10  & 1 & 0 & $-1$ & PE & CEI-2: poder $> 150$ \\
T04 & 50   & $-1$& 1 & 0 & $-1$ & PE & CEI-3: defesa $< 0$ \\
T05 & 50   & 151 & 1 & 0 & $-1$ & PE & CEI-4: defesa $> 150$ \\
T06 & 50   & 10  & 0 & 0 & $-1$ & PE & CEI-5: tipo $< 1$ \\
T07 & 50   & 10  & 4 & 0 & $-1$ & PE & CEI-6: tipo $> 3$ \\
T08 & 50   & 10  & 1 & $-1$ & $-1$ & PE & CEI-7: crítico $< 0$ \\
T09 & 50   & 10  & 1 & 2 & $-1$ & PE & CEI-8: crítico $> 1$ \\
\midrule
% --- Bloco 2: Classes Válidas ---
\multicolumn{8}{l}{\textit{Bloco 2 — Classes Válidas (PE)}} \\
T10 & 100 & 30  & 1 & 0 & 70  & PE & CEV-1: físico normal \\
T11 & 30  & 50  & 1 & 0 & 1   & PE & CEV-2: físico + dano mínimo \\
T12 & 100 & 30  & 1 & 1 & 140 & PE & CEV-3: físico crítico \\
T13 & 30  & 50  & 1 & 1 & 2   & PE & CEV-4: físico crítico + dano mínimo \\
T14 & 100 & 40  & 2 & 0 & 80  & PE & CEV-5: mágico normal \\
T15 & 5   & 12  & 2 & 0 & 1   & PE & CEV-6: mágico + dano mínimo \\
T16 & 100 & 40  & 2 & 1 & 160 & PE & CEV-7: mágico crítico \\
T17 & 5   & 12  & 2 & 1 & 2   & PE & CEV-8: mágico crítico + dano mínimo \\
T18 & 75  & 150 & 3 & 0 & 75  & PE & CEV-9: verdadeiro normal \\
T19 & 75  & 150 & 3 & 1 & 150 & PE & CEV-10: verdadeiro crítico \\
\midrule
% --- Bloco 3: BVA poder ---
\multicolumn{8}{l}{\textit{Bloco 3 — BVA: \texttt{poder\_ataque} {[}1, 150{]}}} \\
T20 & 1   & 0  & 1 & 0 & 1    & BVA & On-min: poder=1 (válido) \\
T21 & 0   & 0  & 1 & 0 & $-1$ & BVA & Off-min: poder=0 (inválido) \\
T22 & 150 & 0  & 1 & 0 & 150  & BVA & On-max: poder=150 (válido) \\
T23 & 151 & 0  & 1 & 0 & $-1$ & BVA & Off-max: poder=151 (inválido) \\
\midrule
% --- Bloco 4: BVA defesa ---
\multicolumn{8}{l}{\textit{Bloco 4 — BVA: \texttt{defesa} {[}0, 150{]}}} \\
T24 & 50  & 0   & 1 & 0 & 50   & BVA & On-min: defesa=0 (válido) \\
T25 & 50  & $-1$& 1 & 0 & $-1$ & BVA & Off-min: defesa=$-1$ (inválido) \\
T26 & 150 & 150 & 1 & 0 & 1    & BVA & On-max: defesa=150 (válido) \\
T27 & 50  & 151 & 1 & 0 & $-1$ & BVA & Off-max: defesa=151 (inválido) \\
\midrule
% --- Bloco 5: BVA tipo ---
\multicolumn{8}{l}{\textit{Bloco 5 — BVA: \texttt{tipo\_ataque} \{1, 2, 3\}}} \\
T28 & 50 & 10 & 1 & 0 & 40   & BVA & On-min: tipo=1 (físico válido) \\
T29 & 50 & 10 & 0 & 0 & $-1$ & BVA & Off-min: tipo=0 (inválido) \\
T30 & 50 & 10 & 3 & 0 & 50   & BVA & On-max: tipo=3 (verdadeiro válido) \\
T31 & 50 & 10 & 4 & 0 & $-1$ & BVA & Off-max: tipo=4 (inválido) \\
\midrule
% --- Bloco 6: BVA crítico ---
\multicolumn{8}{l}{\textit{Bloco 6 — BVA: \texttt{eh\_critico} \{0, 1\}}} \\
T32 & 50 & 10 & 1 & 0  & 40   & BVA & On-min: crítico=0 (normal válido) \\
T33 & 50 & 10 & 1 & $-1$ & $-1$ & BVA & Off-min: crítico=$-1$ (inválido) \\
T34 & 50 & 10 & 1 & 1  & 80   & BVA & On-max: crítico=1 (válido) \\
T35 & 50 & 10 & 1 & 2  & $-1$ & BVA & Off-max: crítico=2 (inválido) \\
\midrule
% --- Bloco 7: BVA dano mínimo ---
\multicolumn{8}{l}{\textit{Bloco 7 — BVA: fronteira do dano mínimo}} \\
T36 & 50 & 50 & 1 & 0 & 1 & BVA & Físico: dano\_base=0 $\to$ mínimo ativado \\
T37 & 51 & 50 & 1 & 0 & 1 & BVA & Físico: dano\_base=1 $\to$ mínimo não ativado \\
T38 & 5  & 10 & 2 & 0 & 1 & BVA & Mágico: dano\_base=0 $\to$ mínimo ativado \\
T39 & 6  & 10 & 2 & 0 & 1 & BVA & Mágico: dano\_base=1 $\to$ mínimo não ativado \\
T40 & 50 & 50 & 1 & 1 & 2 & BVA & Físico+crítico: ordem mín$\to$crítico \\
T41 & 5  & 10 & 2 & 1 & 2 & BVA & Mágico+crítico: ordem mín$\to$crítico \\
\midrule
% --- Bloco 8: Casos adicionais ---
\multicolumn{8}{l}{\textit{Bloco 8 — Casos adicionais}} \\
T42 & 50  & 11  & 2 & 0 & 45  & PE & $\lfloor 11/2 \rfloor = 5$ (divisão inteira ímpar) \\
T43 & 50  & 1   & 2 & 0 & 50  & PE & $\lfloor 1/2 \rfloor = 0$ (defesa=1 não subtrai) \\
T44 & 1   & 150 & 3 & 0 & 1   & PE & Verdadeiro ignora defesa máxima \\
T45 & 150 & 150 & 1 & 1 & 2   & PE & Físico pow=def=150: mín$\to$1$\to$crítico$\to$2 \\
T46 & 150 & 100 & 3 & 1 & 300 & PE & Verdadeiro poder máx + crítico \\
T47 & 50  & 11  & 2 & 1 & 90  & PE & Mágico defesa ímpar + crítico: $45\times2=90$ \\
\end{longtable}

% =========================================================
\section{Entregável 4 — Código-Fonte dos Testes}

\subsection{TestCalcularDano.c}

\begin{lstlisting}[language=C, caption={TestCalcularDano.c --- suíte de 47 testes funcionais.}]
/*
 * TestCalcularDano.c
 * Suite de testes funcionais (caixa-preta) para calcular_dano().
 *
 * Tecnicas aplicadas:
 *   PE  = Particionamento de Equivalencia
 *   BVA = Analise do Valor Limite (Boundary Value Analysis)
 *
 * Parametros e dominios:
 *   poder_ataque : [1, 150]
 *   defesa       : [0, 150]
 *   tipo_ataque  : {1=fisico, 2=magico, 3=verdadeiro}
 *   eh_critico   : {0=normal, 1=critico}
 *
 * Regras (em ordem):
 *   1. Validacao   : parametro fora do dominio => retorna -1
 *   2. Dano base   : fisico=poder-defesa | magico=poder-floor(defesa/2) | verdadeiro=poder
 *   3. Dano minimo : dano < 1 => dano = 1
 *   4. Critico     : eh_critico=1 => dano = dano * 2
 *   5. Retorno     : dano final
 */

#include "unity.h"
#include "unity_fixture.h"
#include "calcular_dano.h"

TEST_GROUP(CalcularDano);
TEST_SETUP(CalcularDano)    { /* nada */ }
TEST_TEAR_DOWN(CalcularDano){ /* nada */ }

/* =========================================================
 * BLOCO 1 - Particionamento de Equivalencia: Classes Invalidas
 * ========================================================= */

/* CEI-1: poder_ataque < 1 (zero) */
TEST(CalcularDano, T01_Invalido_PoderZero)
{
    TEST_ASSERT_EQUAL(-1, calcular_dano(0, 10, 1, 0));
}

/* CEI-1: poder_ataque < 1 (negativo) */
TEST(CalcularDano, T02_Invalido_PoderNegativo)
{
    TEST_ASSERT_EQUAL(-1, calcular_dano(-5, 10, 1, 0));
}

/* CEI-2: poder_ataque > 150 */
TEST(CalcularDano, T03_Invalido_PoderExcede)
{
    TEST_ASSERT_EQUAL(-1, calcular_dano(151, 10, 1, 0));
}

/* CEI-3: defesa < 0 */
TEST(CalcularDano, T04_Invalido_DefesaNegativa)
{
    TEST_ASSERT_EQUAL(-1, calcular_dano(50, -1, 1, 0));
}

/* CEI-4: defesa > 150 */
TEST(CalcularDano, T05_Invalido_DefesaExcede)
{
    TEST_ASSERT_EQUAL(-1, calcular_dano(50, 151, 1, 0));
}

/* CEI-5: tipo_ataque < 1 */
TEST(CalcularDano, T06_Invalido_TipoZero)
{
    TEST_ASSERT_EQUAL(-1, calcular_dano(50, 10, 0, 0));
}

/* CEI-6: tipo_ataque > 3 */
TEST(CalcularDano, T07_Invalido_TipoExcede)
{
    TEST_ASSERT_EQUAL(-1, calcular_dano(50, 10, 4, 0));
}

/* CEI-7: eh_critico < 0 */
TEST(CalcularDano, T08_Invalido_CriticoNegativo)
{
    TEST_ASSERT_EQUAL(-1, calcular_dano(50, 10, 1, -1));
}

/* CEI-8: eh_critico > 1 */
TEST(CalcularDano, T09_Invalido_CriticoExcede)
{
    TEST_ASSERT_EQUAL(-1, calcular_dano(50, 10, 1, 2));
}

/* =========================================================
 * BLOCO 2 - Particionamento de Equivalencia: Classes Validas
 * ========================================================= */

/* CEV-1: Fisico, normal, dano_base > 0 --- dano = 100-30 = 70 */
TEST(CalcularDano, T10_Fisico_Normal)
{
    TEST_ASSERT_EQUAL(70, calcular_dano(100, 30, 1, 0));
}

/* CEV-2: Fisico, normal, minimo ativado --- dano = max(-20,1) = 1 */
TEST(CalcularDano, T11_Fisico_Normal_DanoMinimo)
{
    TEST_ASSERT_EQUAL(1, calcular_dano(30, 50, 1, 0));
}

/* CEV-3: Fisico, critico --- dano = 70*2 = 140 */
TEST(CalcularDano, T12_Fisico_Critico)
{
    TEST_ASSERT_EQUAL(140, calcular_dano(100, 30, 1, 1));
}

/* CEV-4: Fisico, critico, minimo --- dano = min(-20)=1 => 1*2=2 */
TEST(CalcularDano, T13_Fisico_Critico_DanoMinimo)
{
    TEST_ASSERT_EQUAL(2, calcular_dano(30, 50, 1, 1));
}

/* CEV-5: Magico, normal --- dano = 100-floor(40/2)=80 */
TEST(CalcularDano, T14_Magico_Normal)
{
    TEST_ASSERT_EQUAL(80, calcular_dano(100, 40, 2, 0));
}

/* CEV-6: Magico, normal, minimo --- dano = 5-6=-1 => 1 */
TEST(CalcularDano, T15_Magico_Normal_DanoMinimo)
{
    TEST_ASSERT_EQUAL(1, calcular_dano(5, 12, 2, 0));
}

/* CEV-7: Magico, critico --- dano = 80*2=160 */
TEST(CalcularDano, T16_Magico_Critico)
{
    TEST_ASSERT_EQUAL(160, calcular_dano(100, 40, 2, 1));
}

/* CEV-8: Magico, critico, minimo --- dano = -1=>1 => 1*2=2 */
TEST(CalcularDano, T17_Magico_Critico_DanoMinimo)
{
    TEST_ASSERT_EQUAL(2, calcular_dano(5, 12, 2, 1));
}

/* CEV-9: Verdadeiro, normal --- ignora defesa; dano = 75 */
TEST(CalcularDano, T18_Verdadeiro_Normal)
{
    TEST_ASSERT_EQUAL(75, calcular_dano(75, 150, 3, 0));
}

/* CEV-10: Verdadeiro, critico --- dano = 75*2=150 */
TEST(CalcularDano, T19_Verdadeiro_Critico)
{
    TEST_ASSERT_EQUAL(150, calcular_dano(75, 150, 3, 1));
}

/* =========================================================
 * BLOCO 3 - BVA: poder_ataque [1, 150]
 * ========================================================= */

/* BVA poder=1 (on-min): dano = 1-0 = 1 */
TEST(CalcularDano, T20_BVA_PoderMin_On)
{
    TEST_ASSERT_EQUAL(1, calcular_dano(1, 0, 1, 0));
}

/* BVA poder=0 (off-min): invalido */
TEST(CalcularDano, T21_BVA_PoderMin_Off)
{
    TEST_ASSERT_EQUAL(-1, calcular_dano(0, 0, 1, 0));
}

/* BVA poder=150 (on-max): dano = 150-0 = 150 */
TEST(CalcularDano, T22_BVA_PoderMax_On)
{
    TEST_ASSERT_EQUAL(150, calcular_dano(150, 0, 1, 0));
}

/* BVA poder=151 (off-max): invalido */
TEST(CalcularDano, T23_BVA_PoderMax_Off)
{
    TEST_ASSERT_EQUAL(-1, calcular_dano(151, 0, 1, 0));
}

/* =========================================================
 * BLOCO 4 - BVA: defesa [0, 150]
 * ========================================================= */

/* BVA defesa=0 (on-min): dano = 50-0 = 50 */
TEST(CalcularDano, T24_BVA_DefesaMin_On)
{
    TEST_ASSERT_EQUAL(50, calcular_dano(50, 0, 1, 0));
}

/* BVA defesa=-1 (off-min): invalido */
TEST(CalcularDano, T25_BVA_DefesaMin_Off)
{
    TEST_ASSERT_EQUAL(-1, calcular_dano(50, -1, 1, 0));
}

/* BVA defesa=150 (on-max): dano = 150-150=0 => min => 1 */
TEST(CalcularDano, T26_BVA_DefesaMax_On)
{
    TEST_ASSERT_EQUAL(1, calcular_dano(150, 150, 1, 0));
}

/* BVA defesa=151 (off-max): invalido */
TEST(CalcularDano, T27_BVA_DefesaMax_Off)
{
    TEST_ASSERT_EQUAL(-1, calcular_dano(50, 151, 1, 0));
}

/* =========================================================
 * BLOCO 5 - BVA: tipo_ataque {1, 2, 3}
 * ========================================================= */

/* BVA tipo=1 (on-min): fisico valido --- dano = 50-10=40 */
TEST(CalcularDano, T28_BVA_TipoMin_On)
{
    TEST_ASSERT_EQUAL(40, calcular_dano(50, 10, 1, 0));
}

/* BVA tipo=0 (off-min): invalido */
TEST(CalcularDano, T29_BVA_TipoMin_Off)
{
    TEST_ASSERT_EQUAL(-1, calcular_dano(50, 10, 0, 0));
}

/* BVA tipo=3 (on-max): verdadeiro valido --- dano = 50 */
TEST(CalcularDano, T30_BVA_TipoMax_On)
{
    TEST_ASSERT_EQUAL(50, calcular_dano(50, 10, 3, 0));
}

/* BVA tipo=4 (off-max): invalido */
TEST(CalcularDano, T31_BVA_TipoMax_Off)
{
    TEST_ASSERT_EQUAL(-1, calcular_dano(50, 10, 4, 0));
}

/* =========================================================
 * BLOCO 6 - BVA: eh_critico {0, 1}
 * ========================================================= */

/* BVA critico=0 (on-min): normal valido --- dano = 40 */
TEST(CalcularDano, T32_BVA_CriticoMin_On)
{
    TEST_ASSERT_EQUAL(40, calcular_dano(50, 10, 1, 0));
}

/* BVA critico=-1 (off-min): invalido */
TEST(CalcularDano, T33_BVA_CriticoMin_Off)
{
    TEST_ASSERT_EQUAL(-1, calcular_dano(50, 10, 1, -1));
}

/* BVA critico=1 (on-max): critico valido --- dano = 80 */
TEST(CalcularDano, T34_BVA_CriticoMax_On)
{
    TEST_ASSERT_EQUAL(80, calcular_dano(50, 10, 1, 1));
}

/* BVA critico=2 (off-max): invalido */
TEST(CalcularDano, T35_BVA_CriticoMax_Off)
{
    TEST_ASSERT_EQUAL(-1, calcular_dano(50, 10, 1, 2));
}

/* =========================================================
 * BLOCO 7 - BVA: fronteira do dano minimo
 * ========================================================= */

/* Fisico: dano_base=0 (minimo ativado) --- 50-50=0 => 1 */
TEST(CalcularDano, T36_BVA_DanoMin_Fisico_On)
{
    TEST_ASSERT_EQUAL(1, calcular_dano(50, 50, 1, 0));
}

/* Fisico: dano_base=1 (minimo nao ativado) --- 51-50=1 */
TEST(CalcularDano, T37_BVA_DanoMin_Fisico_Off)
{
    TEST_ASSERT_EQUAL(1, calcular_dano(51, 50, 1, 0));
}

/* Magico: dano_base=0 --- 5-floor(10/2)=0 => 1 */
TEST(CalcularDano, T38_BVA_DanoMin_Magico_On)
{
    TEST_ASSERT_EQUAL(1, calcular_dano(5, 10, 2, 0));
}

/* Magico: dano_base=1 --- 6-floor(10/2)=1 */
TEST(CalcularDano, T39_BVA_DanoMin_Magico_Off)
{
    TEST_ASSERT_EQUAL(1, calcular_dano(6, 10, 2, 0));
}

/* Fisico + critico: testa ordem min -> critico
 * 50-50=0 => min=1 => critico => 2
 * (se ordem errada: 0*2=0 => min=1, resultado igual mas caminho errado) */
TEST(CalcularDano, T40_BVA_OrdemMin_Critico_Fisico)
{
    TEST_ASSERT_EQUAL(2, calcular_dano(50, 50, 1, 1));
}

/* Magico + critico: testa ordem min -> critico
 * 5-5=0 => min=1 => critico => 2 */
TEST(CalcularDano, T41_BVA_OrdemMin_Critico_Magico)
{
    TEST_ASSERT_EQUAL(2, calcular_dano(5, 10, 2, 1));
}

/* =========================================================
 * BLOCO 8 - Casos adicionais para detectar defeitos tipicos
 * ========================================================= */

/* Magico com defesa impar: floor(11/2)=5, nao 6 --- dano = 50-5=45 */
TEST(CalcularDano, T42_Magico_DefesaImpar)
{
    TEST_ASSERT_EQUAL(45, calcular_dano(50, 11, 2, 0));
}

/* Magico com defesa=1: floor(1/2)=0 --- dano = 50-0=50 */
TEST(CalcularDano, T43_Magico_DefesaUm)
{
    TEST_ASSERT_EQUAL(50, calcular_dano(50, 1, 2, 0));
}

/* Verdadeiro ignora defesa maxima --- dano = 1 */
TEST(CalcularDano, T44_Verdadeiro_IgnoraDefesaMaxima)
{
    TEST_ASSERT_EQUAL(1, calcular_dano(1, 150, 3, 0));
}

/* Fisico poder=defesa=150, critico --- 150-150=0=>1=>2 */
TEST(CalcularDano, T45_Fisico_MaxMax_Critico)
{
    TEST_ASSERT_EQUAL(2, calcular_dano(150, 150, 1, 1));
}

/* Verdadeiro poder maximo + critico --- 150*2=300 */
TEST(CalcularDano, T46_Verdadeiro_PoderMax_Critico)
{
    TEST_ASSERT_EQUAL(300, calcular_dano(150, 100, 3, 1));
}

/* Magico defesa impar + critico --- 45*2=90 */
TEST(CalcularDano, T47_Magico_DefesaImpar_Critico)
{
    TEST_ASSERT_EQUAL(90, calcular_dano(50, 11, 2, 1));
}
\end{lstlisting}

\subsection{TestCalcularDano\_Runner.c}

\begin{lstlisting}[language=C, caption={TestCalcularDano\_Runner.c --- registra os 47 casos no grupo Unity.}]
#include "unity.h"
#include "unity_fixture.h"

TEST_GROUP_RUNNER(CalcularDano)
{
    /* --- Bloco 1: Classes Invalidas (PE) --- */
    RUN_TEST_CASE(CalcularDano, T01_Invalido_PoderZero);
    RUN_TEST_CASE(CalcularDano, T02_Invalido_PoderNegativo);
    RUN_TEST_CASE(CalcularDano, T03_Invalido_PoderExcede);
    RUN_TEST_CASE(CalcularDano, T04_Invalido_DefesaNegativa);
    RUN_TEST_CASE(CalcularDano, T05_Invalido_DefesaExcede);
    RUN_TEST_CASE(CalcularDano, T06_Invalido_TipoZero);
    RUN_TEST_CASE(CalcularDano, T07_Invalido_TipoExcede);
    RUN_TEST_CASE(CalcularDano, T08_Invalido_CriticoNegativo);
    RUN_TEST_CASE(CalcularDano, T09_Invalido_CriticoExcede);

    /* --- Bloco 2: Classes Validas (PE) --- */
    RUN_TEST_CASE(CalcularDano, T10_Fisico_Normal);
    RUN_TEST_CASE(CalcularDano, T11_Fisico_Normal_DanoMinimo);
    RUN_TEST_CASE(CalcularDano, T12_Fisico_Critico);
    RUN_TEST_CASE(CalcularDano, T13_Fisico_Critico_DanoMinimo);
    RUN_TEST_CASE(CalcularDano, T14_Magico_Normal);
    RUN_TEST_CASE(CalcularDano, T15_Magico_Normal_DanoMinimo);
    RUN_TEST_CASE(CalcularDano, T16_Magico_Critico);
    RUN_TEST_CASE(CalcularDano, T17_Magico_Critico_DanoMinimo);
    RUN_TEST_CASE(CalcularDano, T18_Verdadeiro_Normal);
    RUN_TEST_CASE(CalcularDano, T19_Verdadeiro_Critico);

    /* --- Bloco 3: BVA poder_ataque --- */
    RUN_TEST_CASE(CalcularDano, T20_BVA_PoderMin_On);
    RUN_TEST_CASE(CalcularDano, T21_BVA_PoderMin_Off);
    RUN_TEST_CASE(CalcularDano, T22_BVA_PoderMax_On);
    RUN_TEST_CASE(CalcularDano, T23_BVA_PoderMax_Off);

    /* --- Bloco 4: BVA defesa --- */
    RUN_TEST_CASE(CalcularDano, T24_BVA_DefesaMin_On);
    RUN_TEST_CASE(CalcularDano, T25_BVA_DefesaMin_Off);
    RUN_TEST_CASE(CalcularDano, T26_BVA_DefesaMax_On);
    RUN_TEST_CASE(CalcularDano, T27_BVA_DefesaMax_Off);

    /* --- Bloco 5: BVA tipo_ataque --- */
    RUN_TEST_CASE(CalcularDano, T28_BVA_TipoMin_On);
    RUN_TEST_CASE(CalcularDano, T29_BVA_TipoMin_Off);
    RUN_TEST_CASE(CalcularDano, T30_BVA_TipoMax_On);
    RUN_TEST_CASE(CalcularDano, T31_BVA_TipoMax_Off);

    /* --- Bloco 6: BVA eh_critico --- */
    RUN_TEST_CASE(CalcularDano, T32_BVA_CriticoMin_On);
    RUN_TEST_CASE(CalcularDano, T33_BVA_CriticoMin_Off);
    RUN_TEST_CASE(CalcularDano, T34_BVA_CriticoMax_On);
    RUN_TEST_CASE(CalcularDano, T35_BVA_CriticoMax_Off);

    /* --- Bloco 7: BVA fronteira dano minimo --- */
    RUN_TEST_CASE(CalcularDano, T36_BVA_DanoMin_Fisico_On);
    RUN_TEST_CASE(CalcularDano, T37_BVA_DanoMin_Fisico_Off);
    RUN_TEST_CASE(CalcularDano, T38_BVA_DanoMin_Magico_On);
    RUN_TEST_CASE(CalcularDano, T39_BVA_DanoMin_Magico_Off);
    RUN_TEST_CASE(CalcularDano, T40_BVA_OrdemMin_Critico_Fisico);
    RUN_TEST_CASE(CalcularDano, T41_BVA_OrdemMin_Critico_Magico);

    /* --- Bloco 8: Casos adicionais --- */
    RUN_TEST_CASE(CalcularDano, T42_Magico_DefesaImpar);
    RUN_TEST_CASE(CalcularDano, T43_Magico_DefesaUm);
    RUN_TEST_CASE(CalcularDano, T44_Verdadeiro_IgnoraDefesaMaxima);
    RUN_TEST_CASE(CalcularDano, T45_Fisico_MaxMax_Critico);
    RUN_TEST_CASE(CalcularDano, T46_Verdadeiro_PoderMax_Critico);
    RUN_TEST_CASE(CalcularDano, T47_Magico_DefesaImpar_Critico);
}
\end{lstlisting}

% =========================================================
\section{Entregável 5 — Relatório de Execução e Defeitos Detectados}

Cada versão defeituosa foi testada com a suíte de 47 casos. A Tabela~\ref{tab:defeitos} resume
os testes que falham e o defeito identificado em cada versão.

\begin{table}[H]
\centering
\caption{Defeitos identificados por versão após execução da suíte de testes.}
\label{tab:defeitos}
\begin{tabular}{@{}p{1.3cm}p{2.8cm}p{7.5cm}@{}}
\toprule
\textbf{Versão} & \textbf{Testes que falham} & \textbf{Defeito identificado} \\
\midrule
versao01 & T20, T44 &
    \texttt{poder\_ataque} válido começa em 2. A validação usa \texttt{<= 1}
    (deveria ser \texttt{< 1}), rejeitando \texttt{poder=1} como inválido. \\
\midrule
versao02 & T18, T19, T26, T44, T45 &
    \texttt{defesa} máximo é 149. A validação usa \texttt{>= 150}
    (deveria ser \texttt{> 150}), rejeitando \texttt{defesa=150} como inválido. \\
\midrule
versao03 & T14, T16, T42, T43, T47 &
    Tipo mágico usa \texttt{poder - defesa} em vez de
    \texttt{poder - floor(defesa / 2)} — a divisão por 2 está ausente. \\
\midrule
versao04 & T18, T19, T30, T46 &
    Tipo verdadeiro (3) não ignora a \texttt{defesa}: aplica a fórmula do
    físico (\texttt{poder - defesa}) em vez de retornar apenas \texttt{poder}. \\
\midrule
versao05 & T11, T13, T15, T17, T26, T36, T38, T40, T41, T45 &
    Regra do dano mínimo ausente: a função retorna zero ou valor negativo
    quando \texttt{dano\_base < 1}, em vez de forçar \texttt{dano = 1}. \\
\midrule
versao06 & T12, T13, T16, T17, T19, T34, T40, T41, T45, T46, T47 &
    Multiplicador crítico é $\times 3$ em vez de $\times 2$: usa
    \texttt{dano = dano * 3} na regra 4. \\
\midrule
versao07 & T18, T19, T30, T44, T46 &
    \texttt{tipo\_ataque} máximo é 2. A validação usa \texttt{> 2}
    (deveria ser \texttt{> 3}), rejeitando \texttt{tipo=3} como inválido. \\
\midrule
versao08 & T20, T22, T24 &
    \texttt{defesa} mínimo é 1. A validação usa \texttt{<= 0}
    (deveria ser \texttt{< 0}), rejeitando \texttt{defesa=0} como inválido. \\
\midrule
versao09 & T13, T17, T40, T41, T45 &
    Regras 3 e 4 em ordem invertida: o crítico é aplicado \textit{antes} do
    dano mínimo. Assim \texttt{dano\_base=0} $\times 2 = 0 \to \text{mín} = 1$
    em vez de \texttt{mín=1} $\times 2 = 2$. \\
\midrule
versao10 & T08, T09, T33, T35 &
    Validação de \texttt{eh\_critico} ausente: valores inválidos ($-1$ e $2$)
    não retornam $-1$, sendo tratados como entradas válidas. \\
\bottomrule
\end{tabular}
\end{table}

\subsection{Discussão dos Resultados}

A suíte de 47 casos detectou defeitos em todas as 10 versões fornecidas, sem falsos positivos
na versão de referência (\texttt{correto.o}). Os aspectos mais relevantes observados são:

\begin{itemize}
    \item \textbf{Erros de operador relacional na validação (versões 01, 02, 07, 08):}
    Troca de \texttt{<} por \texttt{<=} (ou \texttt{>} por \texttt{>=}) nos limites de
    validade. Os casos BVA \textit{on-boundary} foram essenciais para detectá-los, pois
    testam exatamente o valor-limite válido que foi erroneamente rejeitado.

    \item \textbf{Fórmula incorreta de dano mágico (versão 03):}
    A ausência da divisão por 2 é detectada pelos casos com defesa par (T14, T16) e
    pelos casos adicionais com defesa ímpar (T42, T43, T47) que distinguem
    $\lfloor\texttt{def}/2\rfloor$ de $\texttt{def}$.

    \item \textbf{Tipo verdadeiro não ignora defesa (versão 04):}
    Os casos T18/T19 (com \texttt{defesa=150}) e T30/T46 revelam o defeito ao
    comparar o resultado esperado quando a defesa deveria ser ignorada.

    \item \textbf{Ausência do dano mínimo (versão 05):}
    Um conjunto amplo de casos com \texttt{dano\_base $\leq 0$} é necessário para cobrir
    esta situação nos três tipos de ataque.

    \item \textbf{Multiplicador incorreto (versão 06):}
    Os casos com crítico onde o resultado esperado é par e exato (ex.: $140$, $160$,
    $300$) detectam facilmente o uso de $\times 3$.

    \item \textbf{Ordem das regras invertida (versão 09):}
    Os casos T40 e T41 foram projetados especificamente para discriminar a ordem
    correta (mínimo antes do crítico): com \texttt{dano\_base=0} e
    \texttt{eh\_critico=1}, o resultado correto é 2 — mas a ordem invertida produziria
    $0 \times 2 = 0 \to \text{mín} = 1$.

    \item \textbf{Ausência de validação do crítico (versão 10):}
    Os casos BVA \textit{off-boundary} T08, T09, T33 e T35 testam diretamente os
    valores inválidos $-1$ e $2$, garantindo que ambos retornem $-1$.
\end{itemize}

% =========================================================
\section*{Conclusão}

O trabalho demonstrou a efetividade das técnicas PE e BVA no projeto de testes funcionais
de caixa-preta. A combinação das duas abordagens gerou uma suíte enxuta (47 casos) com
alta cobertura de defeitos: todos os 10 defeitos injetados foram detectados, com zero
falsos positivos. Os testes de fronteira (BVA) foram determinantes para revelar erros de
operador relacional, enquanto os casos de equivalência válida cobriram as combinações
lógicas das regras de negócio. Os casos adicionais do Bloco 8 provaram ser essenciais para
distinguir defeitos sutis (fórmula mágica, ordem das regras).

\end{document}
